\section{El paradigma de Compound AI Systems}

\subsection{¿Qué es Compound AI Systems?}

\begin{knowledgebox}{Sistema de IA compuesto vs. modelo único}
Un sistema de IA compuesto está formado por múltiples componentes que interactúan entre sí: llamadas a LLM, recuperadores, ejecutores de código, verificadores, etc. Cada componente tiene su propia firma (tipos de entrada/salida) y se combinan mediante lógica de programa. Esto es similar a la «composición de funciones» de la ingeniería de software tradicional.
\end{knowledgebox}

\begin{itemize}
    \item \textbf{RAG} es el sistema compuesto más simple: recuperación + síntesis mediante LLM
    \item \textbf{ReAct} es un sistema compuesto cíclico: iteración de razonamiento y llamadas a herramientas
    \item \textbf{Multi-hop QA}: encadenamiento de múltiples rondas de recuperación y razonamiento
    \item Sistemas más complejos: pipeline de múltiples etapas, que incluye pasos de verificación, reflexión, etc.
\end{itemize}

\subsection{Los problemas del prompting manual}

\begin{warningbox}{Prompting no es sostenible}
Escribir un prompt a mano parece simple, pero en sistemas compuestos provoca:
\begin{itemize}
    \item \textbf{Fragilidad}: cambiar de modelo o modificar el pipeline exige reescribir todos los prompts
    \item \textbf{No composable}: es difícil combinar de forma modular múltiples prompts
    \item \textbf{Difícil de optimizar}: no es posible buscar de forma sistemática el prompt óptimo
    \item \textbf{No portable}: un buen prompt para un modelo puede no funcionar en otro
\end{itemize}
\end{warningbox}

\subsection{Resumen de la sección}

El prompting manual no es sostenible en sistemas compuestos; se necesita un paradigma de programación más abstracto.

